\subsection{Two-Phase NNI Optimization and Model Persistence}

Hyperparameter optimization (HPO) through frameworks like NNI (Neural Network Intelligence) is essential for finding the optimal balance between accuracy and resource consumption on embedded targets. However, NNI is traditionally designed for exploration, where many "trials" are executed to collect metrics, but the final optimized model weights are often not automatically persisted for production deployment.

To bridge this gap, I implemented a \textbf{Two-Phase HPO Protocol} that ensures the best performing model is seamlessly transitioned into the firmware integration stage.

\subsubsection{Phase 1: Exploration and Metric Collection}

In the first phase, NNI manages a large hyperparameter search space. Each trial is executed in a sandbox environment where the model is trained with a specific set of parameters. To minimize disk overhead and maintain a high throughput during exploration, trials do not save the resulting model files (\texttt{.h5}) to disk. Instead, prioritize reporting performance metrics (accuracy, loss, latency) back to the NNI Manager.

\subsubsection{Phase 2: Exploitation and Model Retraining}

Once the exploration budget is exhausted or an optimal configuration is identified, the workflow triggers the exploitation phase. The \texttt{manager.py} script queries the NNI experiment database to identify the hyperparameters of the "Best Trial."

Crucially, instead of manually recreating the environment, the system re-launches the \texttt{trial.py} script one final time with a specific \texttt{RETRAIN\_MODE} environment variable enabled. This triggers a specialized execution path within the trial script:

\begin{itemize}
    \item \textbf{High-Accuracy Training}: The model is trained using the optimal hyperparameters found in Tier 1.
    \item \textbf{Persistence}: Unlike the exploration trials, this final run explicitly saves the model weights to \texttt{best\_model.h5}.
    \item \textbf{Handover}: The saved file is then picked up by the next node in the LangGraph workflow for Edge AI conversion and deployment.
\end{itemize}

This architectural decision solves the "Exploration-Production gap" common in automated HPO toolchains, allowing the agentic workflow to provide a ready-to-flash model as the final output.
